% Introduction

% Motivation
\subsection{Motivation}

%% SPK suggested first few sentences (with modifications):

The interplay between geometry and topology is an important theme in
the study of 3--manifolds.  Since Thurston's work in the late 1970's,
it has been known that any Haken 3--manifold has a geometric
decomposition \cite{thurston}.  Recently Perelman has announced a
proof that any closed, orientable 3--manifold has such a decomposition
(the geometrization conjecture).  However, it is important not only to
know the existence of a geometric structure on a 3--manifold, but also
to understand the relation between that structure and the topological
description of the manifold.

Often a 3--manifold is described topologically, for example by a
combinatorial diagram indicating a knot complement in $S^3$, or a Dehn
filling description (see Definition \ref{def:dehn} below).  However,
relating this topological description to the geometry of the manifold
seems to be very difficult, and remains an important problem in the
area.  There are few tools available to give geometric data from
topological information.  In this paper, we present a new such tool.

We are particularly interested in topological descriptions of
3--manifolds with torus boundary.  These include all knot and link
complements in $S^3$, as well as knot and link complements in other
3--manifolds.  It is a classic result due to Wallace and Lickorish
that every closed, orientable 3--manifold is given by a Dehn filling
on a link complement in $S^3$ (\cite{wallace}, \cite{lickorish}).
Thus these manifolds are of fundamental importance in the subject.

We will also focus only on those manifolds which are hyperbolic.  The
link complements of Wallace and Lickorish above can be taken to be
hyperbolic.  Thurston showed that on any such cusped hyperbolic
manifold ``most'' Dehn fillings are hyperbolic.  That is, only a
finite number of slopes per link component need be excluded, then any
Dehn filling on remaining slopes yields a hyperbolic manifold
\cite{thurston}.  Additionally, for knots in $S^3$, any knot that is
not a torus or satellite knot is known to be hyperbolic
\cite{thurston:bulletin}.  Therefore in some sense, most 3--manifolds
are hyperbolic.  Mostow--Prasad rigidity implies that the hyperbolic
structure on these finite volume manifolds is unique \cite{mostow},
\cite{prasad}.  However, this unique structure is little understood.

Given a hyperbolic manifold with torus boundary components, i.e. given
a \emph{cusped} hyperbolic manifold, the geometric structure on the
manifold allows us to associate a Euclidean similarity structure,
called a cusp shape, to each boundary component.  The cusp shapes of a
manifold are geometric quantities which we will relate to a manifold's
topological description.

There are several reasons one might be interested in knowing the cusp
shape of a manifold.  Cusp shape is an interesting piece of geometric
information by itself.  It also can be used to give a rough lower
bound on the volume of the manifold, by calculating cusp volumes and
using packing arguments such as those of B{\"o}r{\"o}czky
\cite{boroczky}.  (See however \cite{eudave-munoz-luecke}: this lower
bound may not be very sharp.)  More importantly, knowing the cusp
shape, we also obtain information on the geometry of manifolds
obtained by Dehn filling.  For example, we can find all filling slopes
which might not yield a hyperbolic manifold \cite{agol:bounds},
\cite{lackenby:surg}, \cite{hk:univ}.

Our main results give explicit bounds on the change in cusp shape of a
3--manifold with multiple cusps under all but finitely many Dehn
fillings, where one cusp is left unfilled.  This is essentially the
content of Theorems \ref{thm:main-cusp}, \ref{thm:main-norm-length},
and \ref{thm:main-norm-height} below.  The first, Theorem
\ref{thm:main-cusp}, is the most general, bounding cusp information.
In Theorem \ref{thm:main-norm-length}, we give a differential
inequality bounding the normalized lengths of curves on cusps.  In
Theorem \ref{thm:main-norm-height}, we bound the change in height of
the cusp shape.

We also apply these results to the case of knots in $S^3$.  Given only
a diagram of a large class of knots, we are able to determine bounds
on the height of cusps.  This is the content of Theorem
\ref{thm:main-knot-cusp}.  Note that similar results, relating the
geometry of a knot to its diagram, are rare.  In fact, we know only of
results relating volumes to a diagram.  Lackenby found lower bounds on
volumes of alternating knots based on a condition of a diagram, and
upper bounds for general knots \cite{lackenby:alt-volume}.  The lower
bounds were improved by Agol, Storm and Thurston \cite{ast:volumes},
and the upper by Agol and D. Thurston in an appendix to Lackenby's
paper \cite{lackenby:alt-volume}.  Using cone deformations, we found
similar bounds on volumes for another class of knots
\cite{purcell:volume}.  Using negatively curved metrics, we recently
improved these results with Futer and Kalfagianni \cite{fkp:volumes}.

However, as far as we are aware, the results of this paper are the
first of their kind giving actual geometric bounds on cusp shape based
only on a diagram.


\subsection{Cone deformations}

Our primary technique is cone deformation, which we will now describe.
Using properties of cone deformations, we obtain geometric information
on the topological process of Dehn filling.

\begin{define}
  Given a manifold $M$ with torus boundary $\partial M$ and specified
  slope $s$ on $\partial M$, we obtain a new manifold by gluing a
  solid torus to $\partial M$ such that $s$ bounds a disk in the solid
  torus.  This is called the \emph{Dehn filling of $M$ along $s$}, or
  sometimes \emph{Dehn surgery}, and the resulting manifold is denoted
  $M(s)$.  More generally, if $M$ has $n$ torus boundary components
  with slopes $s_1, \dots, s_n$, then we obtain $M(s_1, \dots, s_n)$
  in the same manner.
\label{def:dehn}
\end{define}

A geometric description of Dehn filling was investigated by Thurston
\cite{thurston}.  Given a complete hyperbolic structure on a manifold
$M$ with torus boundary, he showed the structure could be deformed
through incomplete hyperbolic structures with singularities at the
core of the added solid torus to structures whose completion is again
non-singular.  These structures are exactly hyperbolic structures on
Dehn filled manifolds $M(s)$.  The space of deformations is called
hyperbolic Dehn surgery space.

We restrict our attention to particular paths through hyperbolic Dehn
surgery space in which the incomplete hyperbolic manifolds are all
cone manifolds.  That is, the completion gives a manifold with cone
singularities along the core of the added solid torus, with cone angle
$\alpha$.  (A more complete definition of a cone deformation is given
in Definition \ref{def:cone-deform} below.)

We have some control over cone deformations.  Local rigidity results
due to Hodgson and Kerckhoff \cite{hk:cone-rigid}, extended by
Bromberg in the case of geometrically finite manifolds
\cite{bromberg:cone-rigid}, give control on the change in geometry of
the entire manifold, knowing only information on the change in a
neighborhood of the singular locus.

This geometric control has been used recently to prove several
interesting results in 3--manifold theory.  Hodgson and Kerckhoff used
it to give universal bounds on the number of exceptional Dehn surgery
slopes for finite volume manifolds \cite{hk:univ}.  In the infinite
volume case, cone deformations lie at the heart of drilling and
grafting theorems, described by Bromberg in \cite{bromberg:cone}.
These have recently been used to resolve problems in the field of
Kleinian groups, such as the Bers density conjecture
\cite{brock-bromberg:density}, and existence of type preserving
sequences of Kleinian groups \cite{brock-bromberg-evans-souto:tame}.
This paper gives another application of cone deformations, to finite
volume manifolds with cusps.

\subsection{Cusp shapes}

Let $M$ be a manifold admitting a hyperbolic structure, with a torus
boundary component $T$. Under the hyperbolic structure on $M$, $T$
will be a cusp.  We take a horoball neighborhood of this cusp.  Its
boundary is a horospherical torus, which inherits a Euclidean
structure from the hyperbolic structure on $M$.  The structure is
independent of choice of horoball neighborhood, up to scaling of the
Euclidean metric.  Thus we obtain a class of Euclidean similarity
structures on $T$, which we call the \emph{cusp shape} of the torus
$T$.

Now, suppose the hyperbolic structure of $M$ is changing under a cone
deformation with singularities along a link $\Sigma$ in $M$ disjoint
from $T$.  $\Sigma$ is called the singular locus of the deformation.
$T$ will remain a cusp throughout the deformation, and its cusp shape
will be changing with the change of geometry of $M$.  We bound this
change in terms of the change in geometry in a neighborhood of
$\Sigma$.

\begin{theorem}  Suppose we have the following:
  \begin{itemize}
  \item $M$ a complete finite volume hyperbolic 3--manifold with a
    cusp $T$.
  \item $X=X_0$ a complete finite volume hyperbolic 3--manifold which
    can be joined to $M=X_\tau$ by a smooth one--parameter family of
    hyperbolic cone manifolds $X_t$, such that cone angles go from $0$
    at time $t=0$ to $2\pi$ at time $\tau$.
  \item A singular locus, say with $n$ components, such that the tube
    radius about each component is larger than $R_1\geq 0.56$ for all
    $t$.
  \item A horoball neighborhood $U_t$ of the cusp of $X_t$ which
    deforms to $T$ at time $\tau$.  Let $\gamma(t)$ denote the path
    $\partial U_t$ travels through the Teichm{\"u}ller space of the
    torus, endowed with Teichm{\"u}ller metric $\| \cdot \|$.
  \end{itemize}
	Then there exists a parameterization of the deformation such that
  we have the following bound for $\partial U_t$:
  $$2\,\|\gamma'(t)\|^2 \mbox{Area}(\partial U_t) \leq n\,C(R_1)$$
  Here $C(R_1)$ is a strictly decreasing function of $R_1$ which
  approaches $0$ as $R_1$ approaches infinity.
  Moreover, we can assume $\tau \leq (2\pi)^2$.
\label{thm:main-cusp}
\end{theorem}


Also interesting is the application of Theorem \ref{thm:main-cusp}
to lengths of curves on the cusp in question.  Let $\beta$ be a slope
on the cusp torus, and again let $\gamma(t)$ be the path that the cusp
shape $\partial U_t$ traces through the {T}eichm{\"u}ller space of the
torus.  Let $L_t$ be the normalized length of $\beta$ under the
deformation, with derivative $\dot{L} = \frac{d}{dt}L_t$ (see
section \ref{sec:normlength} for more information).  Then $L_t$ will
satisfy:
\begin{equation}
  -\|\gamma'(t)\|\, L_t \leq \dot{L} \leq
   \|\gamma'(t)\|\, L_t.
  \label{eq:norm-length-change}
\end{equation}

Combining equation (\ref{eq:norm-length-change}) with the result of
Theorem \ref{thm:main-cusp}, we may bound the change in normalized
length of a curve under a cone deformation.

\begin{theorem}
Let the set up be as in Theorem \ref{thm:main-cusp} above.  Let $L_t$
be the normalized length of a slope on a cusp torus under the cone
deformation, with derivative $\dot{L} = \frac{d}{dt}L_t$.  Then
%% Aug 6, 2007: note parameterization
there exists a parameterization of the cone deformation such that
%%
the change in $L_t$ is bounded by the following inequality:
$$ -\left(\sqrt{\frac{n\,C(R_1)}{2\mbox{Area}(\partial U_t)}}\right)
\, L_t \leq \dot{L} \leq
\left(\sqrt{\frac{n\,C(R_1)}{2\mbox{Area}(\partial U_t)}}\right) \,
L_t
$$
\label{thm:main-norm-length}
\end{theorem}

Finally, we use Theorem \ref{thm:main-norm-length} to bound the
normalized length of the arc perpendicular to the shortest curve on
the cusp torus, and then bound the length of this shortest curve
below.  We then obtain an estimate that is independent of time.

\begin{theorem}
  Let the notation be as in Theorem \ref{thm:main-cusp} above.  Denote
  by $h(M)$ the normalized \emph{height} of the cusp torus $\partial
  U_\tau$.  That is, $h(M)$ is the normalized length of the arc
  perpendicular to the shortest curve on $\partial U_\tau$ that runs
  from the shortest curve back to the shortest curve (this arc might
  not be closed).  Similarly, let $h(X)$ denote the normalized height
  of $\partial U_0$.  Then
%% Aug 6, 2007:  Parameterization
	there exists a parameterization of the cone deformation such that
%%
	the change in normalized height is bounded in terms of $R_1$ alone:
$$-(2\pi)^2\frac{\sqrt{n\,C(R_1)}}{(1-e^{-2R_1})\sqrt{2}} \leq
 h(M) - h(X) \leq
(2\pi)^2\frac{\sqrt{n\,C(R_1)}}{(1-e^{-2R_1})\sqrt{2}}.$$
\label{thm:main-norm-height}
\end{theorem}

%% Aug 17:  point out that bound on R_1 can be ensured throughout the
%% deformation by selecting slopes with long normalized lengths?

%% Our hyperbolic cone deformation in the above theorems is a hyperbolic
%% Dehn filling along slopes $s_1, \dots, s_n$.  Fix $R_1 \geq 0.56$.
%% Due to work of Hodgson and Kerckhoff \cite[see proof of Theorem
%% 	5.7]{hk:univ}, we can ensure the tube radius about the singular
%% locus is larger than $R_1$ throughout the deformation by selecting the
%% slopes $s_1, \dots, s_n$ to have sufficiently long normalized
%% lengths.  Thus Theorems \ref{thm:main-cusp},
%% \ref{thm:main-norm-length}, and \ref{thm:main-norm-height} could be
%% rewritten in terms of normalized lengths of Dehn filling slopes.  We
%% use these ideas for applications to knot theory.

% Intro to results:  Application to knot theory

\subsection{Applications to knot theory}

In the last part of this paper we give specific applications of
Theorem \ref{thm:main-norm-height}.  We determine bounds on the cusp
shape of a knot based solely on a diagram of that knot.

Before we state the results, we review some definitions.

\begin{define}
  Let $K \subset S^3$ be a knot with diagram $D(K)$.  Define a
  \emph{twist region} to be a region of $D(K)$ where two strands of
  the diagram twist around each other maximally.  Precisely, if $D(K)$
  is viewed as a 4--valent planar graph with over--under crossing
  information at each vertex, then a twist region is a maximal string
  of bigons in the complement of the graph arranged end to end, or a
  single vertex of the graph adjacent to no bigons.  
\end{define}

We also require the diagram to be reduced, in the sense of the
following two definitions.  These definitions are illustrated
elsewhere (see e.g. \cite{lackenby:surg}, \cite{futer-purcell}).

\begin{define}
  A diagram $D(K)$ is \emph{prime} if when any simple closed curve
  $\gamma$ intersects $D(K)$ transversely in two points in the
  interiors of edges, then $\gamma$ bounds a subdiagram containing no
  crossings of the original diagram.
	\label{def:prime}
\end{define}

\begin{define}
  The diagram $D(K)$ is \emph{twist reduced} if when any simple closed
  curve $\gamma$ intersects $D(K)$ transversely in four points in the
  interiors of edges, with two of these points adjacent to one
  crossing and two others adjacent to another crossing, then $\gamma$
  bounds a subdiagram consisting of a (possibly empty) collection of
  bigons arranged end to end between these two crossings.
	\label{def:twred}
\end{define}

Any diagram of a prime knot or link can be simplified into a prime,
twist reduced diagram.  If the diagram is not prime, then crossings on
one side of a simple closed curve are extraneous and can be removed.
If the diagram is not twist reduced, then performing flypes will
amalgamate two twist regions adjacent to a simple closed curve into
one, reducing the number of twist regions.

For a hyperbolic knot complement, we can choose a horoball
neighborhood about the cusp such that a meridian on the boundary of
that horoball has length $1$ (see e.g. \cite{adams:waist}).  Consider
the geodesic arc orthogonal to the meridian which runs from meridian
to meridian on this horospherical torus.  Define the length of this
arc to be the \emph{height} of the cusp.  Note that when the meridian
is the shortest curve on the cusp (as will be the case in the knots of
this paper), this definition of height agrees with that in Theorem
\ref{thm:main-norm-height}.

Let $K$ be a knot in $S^3$ which admits a prime, twist reduced diagram
with at least $2$ twist regions, with each twist region containing at
least $6$ crossings.  Then it was shown in \cite{futer-purcell} that
$S^3-K$ is hyperbolic.  The following theorem gives bounds on the
cusp shape.

\begin{theorem}
  Let $K$ be a knot in $S^3$ which admits a prime, twist reduced
  diagram with a total of $n\geq 2$ twist regions, with each twist
  region containing at least $c\geq 116$ crossings.  In a hyperbolic
  structure on $S^3-K$, take a horoball neighborhood $U$ about $K$.
  Normalize so that the meridian on $\partial U$ has length $1$.  Then
  the height $H$ of the cusp of $S^3-K$ satisfies:
  $$H\geq n\,(1-f(c))^2.$$
  Here $f(c)$ is a positive function of $c$ which approaches $0$ as
  $c$ increases to infinity.

  Additionally, $H\leq n(\sqrt{n-1} + f(c))^2$.
\label{thm:main-knot-cusp}
\end{theorem}

The function $f$ is given explicitly in Section \ref{sec:knots}.  We
can plug in values of $c$ to obtain more specific estimates.  For
example, in Corollary \ref{cor:0.8} we let $c=145$.  Then $1-f(c)$ is
approximately $0.8154$.

Note that the shortest non-meridional slope on the cusp of $S^3-K$
will have length at least that of the height.  So Theorem
\ref{thm:main-knot-cusp} implies that in this case, the shortest
non-meridional slope has length at least $n\,(1-f(c))^2$.

By results of Hodgson and Kerckhoff \cite{hk:univ}, Dehn filling along
a slope with normalized length at least $7.515$ results in a hyperbolic
manifold.  Corollary \ref{cor:0.8} implies that for a complicated
knot, the shortest non-meridional slope will have normalized length at
least $\sqrt{n}(0.8154)$.  Thus the normalized length of a slope will
be larger than $7.515$ whenever $n\geq 85$.  Thus we have the
following corollary to Theorem \ref{thm:main-knot-cusp}:

\begin{corollary}
  Let $K$ be a knot in $S^3$ which admits a prime, twist reduced
  diagram with at least $85$ twist regions and at least $145$
  crossings per twist region.  Then any non-trivial {D}ehn filling of
  $S^3-K$ is hyperbolic.
\label{cor:knot-cusp}
\end{corollary}

Corollary \ref{cor:knot-cusp} is proven without any reference to the
geometrization conjecture.  If we assume that conjecture, then we
recently proved (with Futer \cite{futer-purcell}) that the numbers
$85$ and $145$ could be reduced to $4$ and $6$ respectively.  However,
note the proof of that result gives no geometric information on cusp
geometry.


The numbers of crossings, 116 and 145, can be reduced significantly
for certain classes of knots.  We illustrate this in section
\ref{subsec:2bridge} for 2--bridge knots.


\subsection{Organization of this paper}
\label{sec:organize}

In Section \ref{sec:deform-cone}, we review background information on
cone deformations.

In Section \ref{sec:cusp-est}, we begin the proofs of Theorems
\ref{thm:main-cusp}, \ref{thm:main-norm-length}, and
\ref{thm:main-norm-height} using information set up in Section
\ref{sec:deform-cone}.  This leads us to results which contain certain
boundary terms.

In Section \ref{sec:bdrybound}, we give bounds on these boundary
terms, which allow us to conclude the proof of Theorems
\ref{thm:main-cusp} and \ref{thm:main-norm-length}.

We complete the proof of \ref{thm:main-norm-height} in Section
\ref{sec:merbound}.

Finally, in Section \ref{sec:knots} we apply the results of the
previous sections to knots, and present the proof of Theorem
\ref{thm:main-knot-cusp}.

\subsection{Acknowledgments}

We would like to thank Steve Kerckhoff for many helpful conversations.
We lean heavily on his results with Craig Hodgson.  His willingness to
clarify, explain, and correct has made this paper possible.  We would
also like to thank the referee for very carefully reading this paper,
and for many helpful suggestions, corrections, and comments.

%%%%%%%%%%%%%%%%%%%%%%%%%%%%%%%%%%%%%%%%%%%%%%%%%%%%%%%%%%%%%%%%%


